\documentclass[preprint,12pt]{elsarticle}
\usepackage{amsmath,amssymb,amsthm,amsfonts}
\usepackage{stix}
\usepackage{booktabs,graphicx,xcolor,float,array}
\usepackage[colorlinks=true,linkcolor=blue,citecolor=blue,urlcolor=blue]{hyperref}
\newtheorem{theorem}{Theorem}
\newtheorem{result}{Result}
\theoremstyle{remark}
\newtheorem{remark}{Remark}
\newfloat{algorithm}{tbp}{loa}
\floatname{algorithm}{Algorithm}
\newcommand{\todo}[1]{\textcolor{red}{\textbf{[TODO: #1]}}}

\newcommand{\step}[1]{\par\noindent\hangindent=1.5em\hangafter=1\textbf{Step #1.}\ }
\setlength{\emergencystretch}{2.5em}
\newcommand{\repo}{\url{https://github.com/mossaiby/iHaar}}
\newcommand{\archive}{\url{https://github.com/mossaiby/iHaar/releases/tag/v1.2}}
\journal{Engineering Analysis with Boundary Elements}

\begin{document}
\begin{frontmatter}
\title{Fast solvers on arbitrary grids for Haar wavelet collocation:
sharp-interface phase-field computations and adaptive elastostatics}

\author[inst1]{Farshid Mossaiby\corref{cor1}}
\ead{mossaiby@eng.ui.ac.ir}
\cortext[cor1]{Corresponding author}
\affiliation[inst1]{organization={Department of Civil Engineering, University of Isfahan},
            city={Isfahan}, postcode={81744-73441}, country={Iran}}

\begin{abstract}
The Haar wavelet collocation method (HWCM) is popular in engineering because it is simple to program,
but in two and three dimensions it is usually solved through large dense matrices and is
limited to second-order accuracy. We show that, written in nodal values, the method becomes midpoint
quadratic-spline collocation. An exact three-term relation turns the dense HWCM matrices into tridiagonal
ones, and we prove that the resulting second-derivative operator has real, simple and nonpositive
eigenvalues on every strictly increasing grid. Fast diagonalization, which reduces tensor-grid problems to
a few dense matrix products, is therefore always well defined, including on adaptively refined grids, and
was backward stable on random and adversarial grids. On uniform grids, a cost-free correction raises the
accuracy to fourth order for zero-flux conditions. We present Allen--Cahn and Cahn--Hilliard solvers that
conserve mass exactly, satisfy a modified-energy law, and resolve interfaces with about four nodes across
the transition layer, and a plane elastostatics solver with displacement and traction boundaries, whose
automatic refinement recovered the optimal convergence rate at the tested corner singularities. Against
biquadratic finite elements on identical grids, the method needed about a quarter of the unknowns. All
code is openly available.
\end{abstract}

\begin{keyword}
Haar wavelet collocation \sep spline collocation \sep fast diagonalization \sep graded grids \sep Cahn--Hilliard \sep
sharp interfaces \sep corner singularities \sep adaptive refinement
\end{keyword}
\end{frontmatter}

%==========================================================================
\section{Introduction}
%==========================================================================
The Haar wavelet collocation method (HWCM)~\cite{chen1997haar,lepik2005numerical,lepik2011solving,lepik2014book}
is widely used in applied mechanics for beams, plates, shells and functionally graded
structures~\cite{majak2015accuracy,majak2015convergence,majak2018higher}. The highest derivative of the
unknown is expanded in Haar functions, the lower derivatives follow by exact integration, and the
equation is enforced at cell midpoints. Higher-order variants~\cite{majak2018higher,ratas2021solving} and
graded grids based on non-uniform Haar functions~\cite{dubeau2004nonuniform,lepik2008nonuniform,oruc2019nonuniform}
are available. In two dimensions, however, all equations are usually assembled into one dense matrix of
size $N^2\times N^2$, which becomes impractical beyond about $64\times64$ cells. Graded grids must be
designed by hand, and for nonlinear time-dependent problems no stability results are available.

In terms of nodal values, HWCM is midpoint quadratic-spline collocation
(QSC)~\cite{houstis1988quadratic,christara1994quadratic,christara2006optimal}, and it uses the same
functions as quadratic isogeometric collocation~\cite{auricchio2010isogeometric,schillinger2013isogeometric}.
On tensor grids, separable problems can be solved by fast diagonalization~\cite{lynch1964direct}, which
reduces two- and three-dimensional problems to dense matrix products and fast transforms that run
efficiently on CPUs and GPUs. It requires each one-dimensional operator to be diagonalizable with real
eigenvalues. For
symmetric finite-element and isogeometric Galerkin operators this is automatic~\cite{sangalli2016isogeometric}.
Collocation operators are generally not symmetric, and for them the property has to be proved. For
Hermite-cubic orthogonal spline collocation it is known~\cite{dyksen1987tensor,sun1999spectral}, and it
underlies the matrix decomposition algorithms of Bialecki and Fairweather~\cite{bialecki1995matrix,bialecki2011matrix}.
For QSC, fast solvers exist on uniform grids~\cite{christara2002fast}. To our knowledge, the property has
not been established for midpoint QSC, equivalently nodal HWCM, on arbitrary graded grids. There, the
collocation operator is not symmetric and is not symmetrized by quadrature weights, so a different
argument is needed. Without a guarantee, fast diagonalization on adaptively refined grids could fail.

This paper closes that gap and builds on it. Our central result (Theorem~\ref{thm:spectrum}) is that, on
every grid, the HWCM second-derivative operator has real, distinct and nonpositive eigenvalues, so the
fast solver always applies. We check its numerical reliability (Section~\ref{sec:robust}) and use it in
three engineering settings: fast solution on arbitrary graded grids, Allen--Cahn and Cahn--Hilliard
computations with thin interfaces, which conserve mass exactly, and elastostatics of compressible
materials with traction boundaries and automatic, including local, refinement at corner singularities.
Table~\ref{tab:closest} compares the method with the closest existing approaches. The main text keeps
the mathematics to what is needed to use the method; derivations are in the appendices, and formal
statements and proofs in Part~A of the Supplementary Material. The code is available at \repo.

\begin{table}[t]\centering\footnotesize
\caption{Comparison with the closest spline-collocation approaches, to the best of our knowledge. ``Spectrum'' refers to the one-dimensional collocation operator used by fast diagonalization or matrix decomposition.}\label{tab:closest}
\begin{tabular}{p{0.14\textwidth}p{0.175\textwidth}p{0.13\textwidth}p{0.23\textwidth}p{0.17\textwidth}}
\toprule
Method & Trial space and collocation points & Grids & Spectrum & Fast solver\\
\midrule
Classical HWCM~\cite{lepik2011solving,oruc2019nonuniform} & integrated Haar ($C^1$ quadratics); midpoints & uniform or a~priori graded & not analysed & none (dense $N^2\times N^2$ solve)\\\addlinespace
Hermite-cubic OSC~\cite{dyksen1987tensor,sun1999spectral,bialecki1995matrix} & $C^1$ cubics; two Gauss points per cell & partitions as in the cited works & real, distinct eigenvalues established & matrix decomposition (MDA), FFT on uniform grids\\\addlinespace
QSC~\cite{houstis1988quadratic,christara2002fast,christara2006optimal} & $C^1$ quadratics; midpoints (optimal variants perturbed) & uniform; non-uniform analysed for convergence & symmetric with trigonometric eigenvectors on uniform grids & FFT and MDA on uniform grids\\\addlinespace
Quadratic IGA collocation~\cite{auricchio2010isogeometric,schillinger2013isogeometric,delorenzis2015isogeometric} & B-splines; Greville or superconvergent points & arbitrary knots; local refinement & not established for collocation; Galerkin pencils symmetric~\cite{sangalli2016isogeometric} & direct or Krylov; FDM for the Galerkin counterpart\\\addlinespace
This work & $C^1$ quadratics; midpoints & arbitrary graded tensor grids & real, simple, nonpositive on every grid (Result~\ref{thm:spectrum}) & FDM on any tensor grid, with a safeguard; DCT/DST on uniform grids\\
\bottomrule
\end{tabular}
\end{table}

%==========================================================================
\section{The method}\label{sec:method}
%==========================================================================
\subsection{From Haar wavelets to nodal values}
Let $a=t_0<t_1<\dots<t_N=b$ be the grid, with cells of width $h_k$ and midpoints $m_k$. HWCM expands
$u''$ in Haar functions and integrates twice. On a given grid, the Haar functions span exactly the
functions that are constant on each cell, so the Haar coefficients carry the same information as the
cell values $c_k$ of $u''$. Integrating twice gives piecewise quadratics with continuous slopes, that is,
quadratic splines. We use the values $U_k$ at the cell midpoints, which are the HWCM collocation values,
as unknowns.

The key observation is that the nodal values and the cell values of $u''$ are linked by an exact
three-term relation, derived in~\ref{app:identity}:
\begin{equation}\label{eq:identity}
\mathbf L\mathbf U+\mathbf b=\mathbf A\mathbf c,
\end{equation}
where $\mathbf L$ is a tridiagonal difference matrix, $\mathbf A$ a tridiagonal averaging matrix, and $\mathbf b$ contains the
prescribed end values; $\mathbf U$ and $\mathbf c$ collect the nodal values $U_k$ and the cell values $c_k$. The second derivative at the nodes is therefore $\mathbf c=\mathbf D_2\mathbf U+\mathbf e$, with $\mathbf D_2=\mathbf A^{-1}\mathbf L$
and a boundary vector $\mathbf e$. Instead of the dense integration matrices of classical HWCM, only two
tridiagonal matrices need to be assembled. Slopes follow in closed form as well, which gives the
first-derivative matrix $\mathbf D_1$ and the end slopes needed for traction conditions (Algorithm~1).
Fig.~\ref{fig:pipeline} summarizes how these one-dimensional operators lead to the fast solvers of
Section~\ref{sec:tensor}.

\begin{algorithm}[t]\small\hrule\smallskip
\noindent\textbf{Algorithm 1: one-dimensional operators} (\texttt{ihaar/core.py}, class \texttt{Grid1D}).
Input: knots $t_0<\dots<t_N$ and the type of each end (value prescribed, or zero slope).
\step{1} Compute the cell widths and midpoints, and assemble the tridiagonal matrices $\mathbf L$ and $\mathbf A$ (\ref{app:identity}).
\step{2} Factorize $\mathbf A$ (tridiagonal LU) and form $\mathbf D_2=\mathbf A^{-1}\mathbf L$ and the boundary vectors $\mathbf e_i=\mathbf A^{-1}\mathbf b_i$.
\step{3} Form the first-derivative matrix $\mathbf D_1$ by summing the cell values of $u''$ along the grid~\eqref{eq:D1}.
\step{4} Form the end slopes $s'(a)$ and $s'(b)$ as row vectors, for traction conditions.
\smallskip\hrule
\end{algorithm}

\begin{figure}[t]\centering
\includegraphics[width=0.95\textwidth]{../figures/panels/pipeline.pdf}
\caption{From the grid to the solution. The one-dimensional operators are assembled from two tridiagonal
matrices; Theorem~\ref{thm:spectrum} guarantees a real eigendecomposition on any grid, and the
multidimensional problem is then solved by dense matrix products instead of one dense system.}\label{fig:pipeline}
\end{figure}

\subsection{Main result: the fast solver always works}
\begin{theorem}[Spectrum on every grid]\label{thm:spectrum}
On every strictly increasing grid, and whether a value or a zero slope is prescribed at each end, the
second-derivative matrix $\mathbf D_2=\mathbf A^{-1}\mathbf L$ has real, distinct and nonpositive eigenvalues. It therefore has a
real eigenvector matrix $\mathbf V$, and fast diagonalization is well defined.
\end{theorem}

In plain terms: whatever grid the user or the refinement algorithm produces, the fast solver of
Section~\ref{sec:tensor} can be used, without complex arithmetic and without restrictions such as a bounded
ratio between neighbouring cells. On uniform grids $\mathbf D_2$ is symmetric, and the result is easy. On graded
grids $\mathbf D_2$ is not symmetric, and the proof (Result~2 in the Supplementary Material) uses the theory of
oscillation matrices, developed by Gantmacher and Krein for vibrating mechanical
systems~\cite{gantmacher2002oscillation}. Fig.~\ref{fig:spectrum} illustrates the result on a grid whose
cells differ in size by three orders of magnitude in random order. The theorem says that the solver is
\emph{defined}; whether it is also \emph{accurate} in floating-point arithmetic is checked separately in
Section~\ref{sec:robust}.

\begin{figure}[t]\centering
\begin{tabular}{@{}cc@{}}
\begin{tabular}[t]{@{}c@{}}\raisebox{-\height}{\includegraphics{../figures/panels/spec_grid.pdf}}\end{tabular} &
\begin{tabular}[t]{@{}c@{}}\raisebox{-\height}{\includegraphics{../figures/panels/spec_eigs.pdf}}\end{tabular}\\
\small (a) & \small (b)\\
\end{tabular}
\caption{Theorem~\ref{thm:spectrum} on a strongly non-uniform grid ($N=48$). (a) Cell sizes, spanning three
orders of magnitude in random order. (b) Eigenvalues of $\mathbf D_2$ for both types of end condition: all are
real (to machine precision), distinct and nonpositive.}\label{fig:spectrum}
\end{figure}

\subsection{Uniform grids: cosine transforms, fourth order and mass conservation}
On a uniform grid, $\mathbf D_2$ is symmetric and is diagonalized by the discrete sine transform (prescribed
values) or cosine transform (zero slopes), so the eigenvectors need not be stored. The corrected operator
\begin{equation}\label{eq:D24}
\mathbf D_2^{(4)}=\mathbf D_2+\tfrac{h^2}{24}\mathbf D_2^2
\end{equation}
has the same eigenvectors and is fourth-order accurate for zero-flux and periodic problems; with general
prescribed values, the order remains two. The correction costs nothing, because in the transform domain it
only changes the eigenvalues. Plain HWCM has half the error constant of the three-point difference, so it
is about twice as accurate on the same grid. Finally, with zero slopes at both ends the weighted sum
$\sum_kh_kU_k$ is conserved exactly by $\mathbf D_2$ on any grid, which gives exact mass conservation for the
Cahn--Hilliard equation. Formal statements and proofs of these properties are Results~3 and~4 of the
Supplementary Material.

